%C1===============================[=.............................]===============================%C1
\def\ConeEN{}
\def\ConeLT{Kurios iš žemiau pavaizduotų virvių neįmanoma gauti, lankstant paveikslėlyje dešinėje pavaizduotą virvę? Virvės kirpti negalima.}
\def\ConePL{}
\def\ConeRU{}


\def\rCone{B}
\def\aCone{\granswer[0.4]{C01-1}{C01-2}{C01-3}{C01-4}{C01-5}}


%C2===============================[==............................]===============================%C2
\def\CtwoEN{}
\def\CtwoLT{Kurį paveikslėlį reikia įklijuoti šalia pavaizduotame tuščiame stačiakampyje, kad vaikai laikytų skirtingus aitvarus?}
\def\CtwoPL{}
\def\CtwoRU{}

\def\rCtwo{D}
\def\aCtwo{\granswercadet[0.7]{C02-1}{C02-2}{C02-3}{C02-4}{C02-5}}

%C3===============================[===...........................]===============================%C3
\def\CthreeEN{}
\def\CthreeLT{Paveikslėlyje pavaizduotas rombas ir penkiakampis,  gautas prie rombo pridėjus du stačiuosius trikampius. Kiek procentų  penkiakampio  plotas yra didesnis už rombo plotą?}
\def\CthreePL{}
\def\CthreeRU{}

\def\rCthree{E}
\def\aCthree{\answer{20\%}{25\%}{30\%}{40\%}{50\%}}

%C4===============================[====..........................]===============================%C4
\def\CfourEN{}
\def\CfourLT{${\dfrac{20 \times 24}{2 \cdot 0+2 \cdot 4}}=$}
\def\CfourPL{}
\def\CfourRU{}

\def\rCfour{D}
\def\aCfour{\answer{12}{30}{48}{60}{120}}


%C5===============================[=====.........................]===============================%C5
\def\CfiveEN{}
\def\CfiveLT{Adomas nupjovė keturis tetraedro  kampus, kaip parodyta paveikslėlyje. Kiek viršūnių turi gautas briaunainis?}
\def\CfivePL{}
\def\CfiveRU{}

\def\rCfive{D}
\def\aCfive{\answer{8}{9}{11}{12}{15}}

%C6===============================[======........................]===============================%C6
\def\CsixEN{}
\def\CsixLT{Ema turi tris žetonus, ant kurių parašyti skaičiai 1, 5 ir 11 (žr.~pav.). Ji nori juos sudėti paeiliui, kad susidarytų keturženklis skaičius. Kiek skirtingų keturženklių skaičių Ema gali gauti tokiu būdu?}
\def\CsixPL{}
\def\CsixRU{}

\def\rCsix{B}
\def\aCsix{\answer{3}{4}{6}{8}{9}}

%C7===============================[=======.......................]===============================%C7
\def\CsevenEN{}
\def\CsevenLT{Ant stalo gulėjo penki skanėstai: 
\includegraphics[scale=0.5]{C07-1},
\includegraphics[scale=0.5]{C07-2},
\includegraphics[scale=0.5]{C07-3},
\includegraphics[scale=0.5]{C07-4} ir
\includegraphics[scale=0.5]{C07-5}.
Elena mėgsta \includegraphics[scale=0.5]{C07-1}. Ema mėgsta \includegraphics[scale=0.5]{C07-1},
    \includegraphics[scale=0.5]{C07-3}, \includegraphics[scale=0.5]{C07-4} ir \includegraphics[scale=0.5]{C07-5}.
Austėja mėgsta \includegraphics[scale=0.5]{C07-2}, \includegraphics[scale=0.5]{C07-3},
    \includegraphics[scale=0.5]{C07-4} ir \includegraphics[scale=0.5]{C07-5}. Gerda mėgsta
    \includegraphics[scale=0.5]{C07-1}, \includegraphics[scale=0.5]{C07-2} ir \includegraphics[scale=0.5]{C07-3}.
Adomas mėgsta \includegraphics[scale=0.5]{C07-1} ir \includegraphics[scale=0.5]{C07-3}.
Kiekvienas vaikas gavo tokį skanėstą, kokį mėgsta. Kas gavo \includegraphics[scale=0.5]{C07-3}?}
\def\CsevenPL{}
\def\CsevenRU{}

\def\rCseven{E}
\def\aCseven{\answer{Elena}{Ema}{Austėja}{Gerda}{Adomas}}

%C8===============================[========......................]===============================%C8
\def\CeightEN{}
\def\CeightLT{Statinės tūris yra lygus bendram 12 vienodų didelių butelių tūriui. Jis taip pat lygus bendram 20 vienodų mažų butelių tūriui. Į tuščią statinę supiltas 9 didelių pilnų butelių turinys. Kiek daugiausiai mažų pilnų butelių turinys dar tilps į šią statinę?}
\def\CeightPL{}
\def\CeightRU{}

\def\rCeight{C}
\def\aCeight{\answer{3}{4}{5}{6}{8}}

%C9===============================[=========.....................]===============================%C9
\def\CnineEN{}
\def\CnineLT{Į kiekvieną lentelės langelį Agota įrašė po vieną  natūralųjį skaičių, o paskui juos uždengė. Keturi įrašytieji skaičiai yra skirtingi. Kiekvienos eilutės ir kiekvieno stulpelio skaičių sandaugos nurodytos paveikslėlyje. Kam lygi keturių lentelėje įrašytųjų skaičių suma?}
\def\CninePL{}
\def\CnineRU{}

\def\rCnine{C}
\def\aCnine{\answer{10}{12}{13}{14}{15}}

%C10===============================[==========....................]===============================%C10
\def\CtenEN{}
\def\CtenLT{Prekybos centre keturių vienodų vežimėlių, maksimaliai įstumtų vienas į kitą, ilgis yra 108 cm. 
Dešimties tokių vežimėlių, maksimaliai įstumtų vienas į kitą, ilgis yra 168 cm. Kam lygus vieno vežimėlio ilgis centimetrais?}
\def\CtenPL{}
\def\CtenRU{}

\def\rCten{C}
\def\aCten{\answer{60}{68}{78}{88}{90}}

%C11===============================[===========...................]===============================%C11
\def\CuoneEN{}
\def\CuoneLT{Kotryna iškepė pyragą ir supjaustė jį į dešimt vienodų dalių (tiesiomis linijomis nuo pyrago centro). Ji suvalgė vieną dalį, o likusias devynias išdėliojo vienodais tarpais, kaip parodyta paveikslėlyje. Koks yra kampas tarp gretimų pyrago dalių?}
\def\CuonePL{}
\def\CuoneRU{}

\def\rCuone{B}
\def\aCuone{\answer{$5^\circ$}{$4^\circ$}{$3^\circ$}{$2^\circ$}{$1^\circ$}}

%C12===============================[============..................]===============================%C12
\def\CutwoEN{}
\def\CutwoLT{Evelina nori sudėti $4\times 4$ lentelę, kurioje kiekvienos eilutės ir kiekvieno stulpelio keturių skaičių sumos būtų vienodos. 
Ji turi tris lentelės dalis \includegraphics[scale=0.6]{C12-1}, ir reikia dar vienos. Kurią dalį Evelina turi panaudoti, kad sudėtų tokią lentelę?}
\def\CutwoPL{}
\def\CutwoRU{}

\def\rCutwo{A}
\def\aCutwo{\granswer[0.6]{C12-2}{C12-3}{C12-4}{C12-5}{C12-6}}


%C13===============================[=============.................]===============================%C13
\def\CuthreeEN{}
\def\CuthreeLT{Kvadrate, kurio kraštinės ilgis lygus 10, nubrėžtos trys atkarpos.  
Dviejų užtušuotų trikampių plotai lygūs atitinkamai $A$ ir $B$, kaip parodyta paveikslėlyje. Kam lygus skirtumas $A-B$?}
\def\CuthreePL{}
\def\CuthreeRU{}

\def\rCuthree{A}
\def\aCuthree{\answer{0}{1}{2}{5}{10}}

%C14===============================[==============................]===============================%C14
\def\CufourEN{}
\def\CufourLT{Pingvinė Pipa kasdien eina žvejoti ir visada parneša 12 žuvų dviem savo pingviniukams. 
Kiekvieną dieną pirmajam sutiktam pingviniukui ji duoda 7 žuvis, o antrajam -- 5 žuvis. Pingviniukai iš karto  suėda visas gautas žuvis. Per kelias pastarąsias dienas vienas pingviniukas suėdė 44 žuvis. Kiek žuvų per tas dienas suėdė kitas pingviniukas?}
\def\CufourPL{}
\def\CufourRU{}

\def\rCufour{D}
\def\aCufour{\answer{34}{40}{46}{52}{58}}

%C15===============================[===============...............]===============================%C15
\def\CufiveEN{}
\def\CufiveLT{Dominykas turi daug vienodų kubelių. Dominykas prie vieno kubelio kiekvienos sienos priklijavo po kubelį (žr. pav.). Prie gautojo briaunainio kiekvienos sienos jis nori priklijuoti dar po vieną kubelį. Kiek dar kubelių prireiks Dominykui?}
\def\CufivePL{}
\def\CufiveRU{}

\def\rCufive{A}
\def\aCufive{\answer{18}{16}{14}{12}{10}}

%C16===============================[================..............]===============================%C16
\def\CusixEN{}
\def\CusixLT{Kengūra nušuoliavo į kalno viršūnę, o tada tuo pačiu keliu nušuoliavo žemyn ir grįžo į tą pačią vietą. 
Kiekvienas kengūros šuolis kelyje į kalno viršūnę yra 1 metro ilgio. Kiekvienas jos šuolis žemyn nuo kalno yra tris kartus ilgesnis už kiekvieną šuolį į kalno viršūnę. Kengūra iš viso atliko 2024 šuolius. Kokį atstumą metrais ji įveikė?}
\def\CusixPL{}
\def\CusixRU{}

\def\rCusix{D}
\def\aCusix{\answer{506}{1012}{2024}{3036}{4048}}

%C17===============================[=================.............]===============================%C17
\def\CusevenEN{}
\def\CusevenLT{Sofija padalijo stačiakampį į keturis mažesnius. Trijų gautųjų stačiakampių perimetrai yra 16, 18 ir 24, kaip parodyta paveikslėlyje. Kam lygus ketvirtojo stačiakampio perimetras?}
\def\CusevenPL{}
\def\CusevenRU{}

\def\rCuseven{B}
\def\aCuseven{\answer{8}{10}{12}{14}{16}}

%C18===============================[==================............]===============================%C18
\def\CueightEN{}
\def\CueightLT{Vanduo sudaro 80\% šviežių grybų (pagal masę). Vanduo taip pat sudaro 20\% džiovintų grybų (pagal masę).
Kiek procentų sumažėja grybų masė džiovinimo metu?}
\def\CueightPL{}
\def\CueightRU{}

\def\rCueight{C}
\def\aCueight{\answer{60}{70}{75}{80}{85}}

%C19===============================[===================...........]===============================%C19
\def\CunineEN{}
\def\CunineLT{Kiek skaitmenų turi mažiausias natūralusis skaičius, kurio skaitmenų sandauga yra lygi $2^{2024}$?}
\def\CuninePL{}
\def\CunineRU{}

\def\rCunine{C}
\def\aCunine{\answer{506}{674}{675}{1012}{2024}}

%C20===============================[====================..........]===============================%C20
\def\CutenEN{}
\def\CutenLT{Ant stalo gulėjo devynios kortelės. Kiekvienoje kortelėje užrašytas vienas iš natūraliųjų skaičių nuo 1 iki 9. Užrašytieji skaičiai yra skirtingi. Elena, Elžbieta, Gerda ir Simonas paėmė po dvi korteles. Paaiškėjo, kad Elenos skaičių suma lygi 6, Elžbietos  skaičių skirtumas lygus 5, Gerdos skaičių sandauga lygi 18, o Simono vienas skaičius yra dvigubai didesnis už kitą jo skaičių. Kuris skaičius liko ant stalo?}
\def\CutenPL{}
\def\CutenRU{}

\def\rCuten{E}
\def\aCuten{\answer{1}{3}{6}{8}{9}}

%C21===============================[=====================.........]===============================%C21
\def\CdoneEN{}
\def\CdoneLT{Skaitmenis nuo 0 iki 9 galima sudaryti iš horizontalių ir vertikalių brūkšnelių, kaip parodyta paveikslėlyje.\\
\centerline{\includegraphics[scale=0.3]{C21-1}}\\
Elzė pasirinko tris skirtingus tokius skaitmenis. Šiuos skaitmenis iš viso sudaro 5 horizontalūs ir 10 vertikalių brūkšnelių. Kam lygi Elzės išsirinktų skaitmenų suma?}
\def\CdonePL{}
\def\CdoneRU{}

\def\rCdone{A}
\def\aCdone{\answer{9}{10}{14}{18}{19}}

%C22===============================[======================........]===============================%C22
\def\CdtwoEN{}
\def\CdtwoLT{Didžiausias pirminis skaičius, mažesnis už 50 ir lygus dviejų pirminių skaičių sumai, yra}
\def\CdtwoPL{}
\def\CdtwoRU{}

\def\rCdtwo{D}
\def\aCdtwo{\answer{37}{39}{41}{43}{47}}

%C23===============================[=======================.......]===============================%C23
\def\CdthreeEN{}
\def\CdthreeLT{Į stačiakampį įbrėžti trys pusapskritimiai, kaip parodyta paveikslėlyje. Kam lygus šio stačiakampio perimetras centimetrais?}
\def\CdthreePL{}
\def\CdthreeRU{}

\def\rCdthree{B}
\def\aCdthree{\answer{82}{92}{96}{108}{120}}

%C24===============================[========================......]===============================%C24
\def\CdfourEN{}
\def\CdfourLT{Kiek yra keturženklių natūraliųjų skaičių, kurie dalijasi iš 8, 10 ir 12, bet nesidalija iš~9?}
\def\CdfourPL{}
\def\CdfourRU{}

\def\rCdfour{C}
\def\aCdfour{\answer{6}{8}{50}{52}{66}}

%C25===============================[=========================.....]===============================%C25
\def\CdfiveEN{}
\def\CdfiveLT{Elena ketina į lentelės kiekvieną langelį taip įrašyti po vieną natūralųjį skaičių, kad kiekviename vidurinės ir viršutinės eilutės langelyje įrašytas skaičius būtų lygus po juo esančiuose dviejuose gretimuose langeliuose įrašytų skaičių sandaugai. Ji jau įrašė skaičių 720 (žr. pav.). Keliais būdais Elena gali įrašyti skaičių $n$?}
\def\CdfivePL{}
\def\CdfiveRU{}

\def\rCdfive{B}
\def\aCdfive{\answer{8}{6}{5}{4}{1}}

%C26===============================[==========================....]===============================%C26
\def\CdsixEN{}
\def\CdsixLT{Raminta lentoje užrašė skaičius 4, 6, 12, 13, 22 ir 29 ir kiekvieną iš jų nuspalvino mėlynai arba geltonai. Julija vieną iš šių skaičių nutrynė. Paaiškėjo, kad lentoje likusių geltonųjų skaičių suma yra du kartus didesnė už lentoje likusių mėlynųjų skaičių sumą. Kurį skaičių Julija nutrynė?}
\def\CdsixPL{}
\def\CdsixRU{}

\def\rCdsix{E}
\def\aCdsix{\answer{4}{12}{13}{22}{29}}

%C27===============================[===========================...]===============================%C27
\def\CdsevenEN{}
\def\CdsevenLT{Paveikslėlyje pavaizduoti trys kampai $\alpha$, $\beta$ ir $\gamma$. 
Kam lygi suma $\alpha + \beta + \gamma$?}
\def\CdsevenPL{}
\def\CdsevenRU{}

\def\rCdseven{D}
\def\aCdseven{\answer{$60^\circ$}{$70^\circ$}{$75^\circ$}{$90^\circ$}{$120^\circ$}}

%C28===============================[============================..]===============================%C28
\def\CdeightEN{}
\def\CdeightLT{Kapitonas Flintas paprašė keturių piratų pasakyti, kiek aukso, sidabro ir vario monetų yra slėptuvėje. Jų atsakymai buvo užrašyti popieriaus lape, kuris, deja, buvo apgadintas (žr. pav.). Tik vienas piratas pasakė tiesą, kiti melavo, kiekvieną monetų skaičių nurodydami neteisingai. Kuris piratas sakė tiesą, jei žinoma, kad slėptuvėje iš viso buvo 30 monetų?}
\def\CdeightPL{}
\def\CdeightRU{}

\def\rCdeight{B}
\def\aCdeight{\answer{Tomas}{Simas}{Jonas}{Lukas}{Neįmanoma nustatyti}}

%C29===============================[=============================.]===============================%C29
\def\CdnineEN{}
\def\CdnineLT{Elzė važiuoja iš taško A į tašką B ir iš karto grįžta į tašką A. Gerda važiuoja iš taško B į tašką A ir iš karto grįžta į tašką B. Abi jos važiuoja tuo pačiu keliu, pradeda važiuoti tuo pačiu metu ir kiekviena abiem kryptimis važiuoja pastoviu greičiu. Elzės greitis yra tris kartus didesnis už Gerdos greitį. Pirmą kartą jos susitinka praėjus 15 minučių po starto. Po kiek minučių nuo starto jos susitiks antrą kartą?}
\def\CdninePL{}
\def\CdnineRU{}

\def\rCdnine{C}
\def\aCdnine{\answer{20}{25}{30}{35}{45}}


%C30===============================[==============================]===============================%C30
\def\CdtenEN{}
\def\CdtenLT{Paveikslėlyje pavaizduotas penkiakampis $ABCDE$, kuriame ${\angle A =\angle B = 90^\circ}$, $AE = BC$ ir $ED = DC$.
Atkarpoje $AB$ pažymėti keturi taškai, dalijantys šią atkarpą į penkias lygias dalis. Iš šių keturių taškų nubrėžtos
atkarpos, statmenos kraštinei $AB$ (žr.~pav.). Juodosios dalies plotas lygus 13, o pilkosios dalies plotas lygus 10. Kam lygus penkiakampio $ABCDE$ plotas?}
\def\CdtenPL{}
\def\CdtenRU{}

\def\rCdten{A}
\def\aCdten{\answer{45}{47}{49}{58}{60}}


